\subsection{Small disjoint stalks at larger centers}

\begin{lemma}[Proved: two binary core partitions]
\label{CP2:A-star-two-binary-band-cores}
Let $M$ have a rectangular-band minimal ideal $K=I\times J$.
A source star with two trivial leaves, both stalks binary,
cannot have the rational homology and cohomology ring of $\CP^2$.
\end{lemma}
\begin{proof}
Write $S,T$ for the two right acts. A global $M$-fixed point
in either stalk makes every global second class restrict to
zero on $BM$: the fixed-point section of that stalk's action
projection factors through the star bridge, and the leaf
corner is trivial. The other leaf corner is trivial as well,
so the two-block relative-cup lemma kills every second square.
Assume that neither act has a fixed point.

A singleton core orbit is globally fixed. Indeed if $xk=x$
for every $k\in K$, then $xm=x(km)=x$ for every $m\in M$.
The two binary acts therefore each consist of one two-point
core orbit. For either such orbit, the right action of
$e_{ij}$ is constant with value depending only on $j$.
Thus there are surjections
\[
 \pi_S:J\to S,\qquad \pi_T:J\to T.
\]
For the column actor $\tau_m$ of $m\in M$ they obey
$\pi_S(j)m=\pi_S(\tau_m(j))$ and the corresponding identity
for $T$.

If the two induced partitions of $J$ coincide, the bijection
$\pi_S(j)\mapsto\pi_T(j)$ is well-defined and $M$-equivariant.
Lemma~\ref{CP2:A-star-general-leaf-merge} then transfers a
rational projective-plane target to the ordinary binary act
cone on $T$. The universal binary-act product obstruction
excludes that target.

Suppose the partitions differ. Form a bipartite graph with
vertex set $S\sqcup T$ and one edge
$\pi_S(j)\!-!\pi_T(j)$ for each $j\in J$. Every vertex is
incident to an edge. If this graph were disconnected, its
four vertices would split into two components, each containing
one vertex of $S$ and one of $T$. This would make the two
partitions equal. Hence the graph is connected.

Identify the two trivial leaf vertices. By
Lemma~\ref{CP2:A-star-leaf-reductions}(3), the resulting ordinary
act cone on $S\sqcup T$ has the homotopy type of the original
star wedged with a circle; the quotient induces an isomorphism
on cohomology in degrees at least two. The constant-action
quotient identifies $\pi_S(j)$ with $\pi_T(j)$ for every
column $j$. The connected graph just constructed makes its
state set a singleton. Lemma~\ref{CP2:A-constant-quotient-H2}
therefore gives a surjection from zero second cohomology to
that of the disjoint-union act cone. This contradicts the
required nonzero second group. No connectedness of the original
action category, or vanishing of its first cone homology,
has been assumed.
\end{proof}

\begin{theorem}[Proved: two binary stalks through center order thirteen]
\label{CP2:A-star-thirteen-two-binary-stalks}
Let a source star have center monoid of order at most thirteen,
both leaf corner orders at most two, and nonempty stalks of
size at most two. It cannot have the rational homology and
cohomology ring of $\CP^2$. The opposite assertion also holds.
\end{theorem}
\begin{proof}
The trivial right rational modules of both leaf corners are
projective. The actual orbit functors of
Lemma~\ref{CP2:A-star-leaf-reductions}(1) replace the leaves by
trivial corners and the stalks by nonempty sets of at most two
states, preserving rational homology and its cohomology ring.
A singleton or a fixed state in either stalk kills all squares
by the fixed-point argument of the preceding lemma. We may
assume both stalks are binary and fixed-point-free.

Identify the trivial leaves and put $X$ for the resulting
ordinary cone on their disjoint union. This replaces the star
by its wedge with a circle, so $X$ has bounded rational
homology and retains its nonzero second square. The relative
and two-row act obstructions require nonzero corner second
cohomology, at least three minimal rows, and at least two
minimal columns. Write $r,c,g$ for the row, column and local
group orders of the minimal ideal. If $g\geq2$, then
$rcg\geq12$, and the bound $|M|\leq13$ leaves at most one
outside element. Lemma~\ref{CP2:A-two-outside-acts} would make
$H_4(X;\mathbb Q)=0$, a contradiction. Thus $g=1$ and the
minimal ideal is a rectangular band. The preceding lemma
excludes precisely the remaining star.
\end{proof}

\begin{theorem}[Proved: binary and ternary trivial leaves through order twelve]
\label{CP2:A-star-twelve-binary-ternary-stalks}
Let a source star have trivial leaf corners, center monoid of
order at most twelve, one binary stalk, and another nonempty
stalk of size at most three. It cannot have the rational
homology and cohomology ring of $\CP^2$.
The opposite assertion also holds.
\end{theorem}
\begin{proof}
Write $S$ for the binary stalk and $T$ for the other stalk.
If either contains a fixed point, all squares vanish by the
same fixed-point argument. If $|T|\leq2$, apply the preceding
theorem. Assume $|T|=3$ and neither stalk has a fixed point.

The cone on $S\sqcup T$, obtained by identifying the leaves,
has bounded rational homology and a nonzero second square.
As above, its minimal ideal has at least three rows and two
columns. A nontrivial local group would require at least
thirteen monoid elements, so the ideal is a rectangular band.
If it has two columns, a minimal right ideal has two elements.
Lemma~\ref{CP2:A-binary-act-two-element-right-ideal} makes
$B(S,M)$ contractible. Deleting the binary edge is therefore
an actual homotopy equivalence onto the ordinary act cone on
$T$. This would have the rational projective-plane target,
excluded by Proposition~\ref{CP2:A-twelve-three-state-rational}.

With at least three columns and rows, the order-twelve budget
forces a $3\times3$ minimal band. The binary stalk $S$ is one
two-point core orbit, with column partition $\pi_S$.
The core of $T$ has no singleton orbit, so it consists of
one orbit of size two or three. If it has size three,
then $T=TK$ is the full three-point column orbit. The column
map $T\to S$ is $M$-equivariant, and the general leaf-merge
lemma transfers the rational target to a binary cone, impossible.

In the remaining case write $T=U\sqcup\{w\}$, where $U$
is a two-point core orbit. If its column partition agrees
with $\pi_S$, the labelled column map gives an equivariant
bijection $S\to U\subset T$. Merging the $S$ leaf into the
$T$ leaf transfers the rational target to a three-state act
cone over $M$, again excluded by the rational order-twelve
theorem.

If the two column partitions differ, their bipartite graph
is connected, by the first lemma. In the constant-action
quotient of $S\sqcup T$, all four core points are consequently
identified. Every minimal-band image lies in that core, so
the sole transient point $w$ stays separate and the quotient
has exactly two states. Its degree-two products vanish by
the universal binary-act obstruction. The natural second
cohomology surjection in
Lemma~\ref{CP2:A-constant-quotient-H2} kills all products on
the original disjoint-union cone, and hence on the star.
This contradicts its rational projective-plane ring.
\end{proof}

\begin{corollary}[Proved: the three large-center star profiles]
\label{CP2:A-nineteen-three-large-center-stars}
The nineteen-arrow source-star profiles
\[
 (12,1,1;2,3),\qquad (12,1,2;2,2),\qquad (13,1,1;2,2)
\]
are excluded. The corresponding sink stars are excluded too.
\end{corollary}
\begin{proof}
Use the binary--ternary theorem for the first profile and
the two-binary theorem for the last two.
\end{proof}
